% The references of the paper, shared by main.tex and supplement.tex so that each work has the same number in both;
% main.tex records how the entries were checked.
\bibitem{linnik44} U. V. Linnik, \emph{On the least prime in an arithmetic progression. I. The basic theorem},
  Rec. Math. [Mat. Sbornik] N.S. 15(57) (1944), 139--178.
\bibitem{linnik44b} U. V. Linnik, \emph{On the least prime in an arithmetic progression. II. The Deuring--Heilbronn
  phenomenon}, Rec. Math. [Mat. Sbornik] N.S. 15(57) (1944), 347--368.
\bibitem{heath-brown92} D. R. Heath-Brown, \emph{Zero-free regions for Dirichlet $L$-functions, and the least prime
  in an arithmetic progression}, Proc. London Math. Soc. (3) 64 (1992), 265--338.
\bibitem{xylouris-acta} T. Xylouris, \emph{On the least prime in an arithmetic progression and estimates for the zeros of
  Dirichlet $L$-functions}, Acta Arith. 150 (2011), 65--91.
\bibitem{xylouris11} T. Xylouris, \emph{\"Uber die Nullstellen der Dirichletschen $L$-Funktionen und die kleinste
  Primzahl in einer arithmetischen Progression}, Bonner Math. Schriften 404, Bonn, 2011.
\bibitem{xylouris18} T. Xylouris, \emph{Linniks Konstante ist kleiner als 5}, Chebyshevskii Sb. 19 (2018), 80--94.
\bibitem{lls} Y. Lamzouri, X. Li and K. Soundararajan, \emph{Conditional bounds for the least quadratic non-residue
  and related problems}, Math. Comp. 84 (2015), 2391--2412.
\bibitem{lls-corr} Y. Lamzouri, X. Li and K. Soundararajan, \emph{Corrigendum to ``Conditional bounds for the least
  quadratic non-residue and related problems''}, Math. Comp. 86 (2017), 2551--2554 (it corrects Lemma~6.1 and
  Theorems~1.2 and~1.3 of \cite{lls}).
\bibitem{openai-main} OpenAI, \emph{The Quasi-Riemann Hypothesis: A Zero-Free Half-Plane $\operatorname{Re}(s) > 7/8$},
  preprint, 30 September 2026, \url{https://github.com/openai/math}, entry 003, read at commit \texttt{adc7f12}. Not
  refereed.
\bibitem{openai-companion} OpenAI, \emph{The Quasi-Riemann Hypothesis}, preprint, 5 October 2026, which proves
  $\operatorname{Re}(s) > 11/12$ by a different route from \cite{openai-main}; \url{https://github.com/openai/math},
  entry 003, file
  \file{preprints/The-Quasi-Riemann-Hypothesis-October-5-2026/build/paper2.tex}, read at commit \texttt{adc7f12}.
  Not refereed.
\bibitem{openai-lean} OpenAI, \emph{Catalog of papers with a formalized main result}, file
  \texttt{lean/formalization.yaml} (version v0.4) of the repository \texttt{openai/math}, a web page,
  \url{https://github.com/openai/math/blob/main/lean/formalization.yaml}, read 25 October 2026; with
  \texttt{lean/OAI/NumberTheory/DirichletL/Nonvanishing.lean} of the same repository. It is OpenAI's own record; it
  marks its scope as partial progress and its review status as unchecked.
\bibitem{cgl} B. Chen, V. Gupta and Y. C. Li, \emph{Large Value Estimates for Dirichlet Polynomials with Characters
  and Zero Density of Dirichlet $L$-Functions}, arXiv:2507.08296v2 (preprint; v2 of 27 July 2026, v1 of 11 July 2025
  was by B. Chen alone).
\bibitem{guth-maynard} L. Guth and J. Maynard, \emph{New large value estimates for Dirichlet polynomials},
  Ann. of Math. (2) 203 (2026), 623--675; arXiv:2405.20552.
\bibitem{hb79b} D. R. Heath-Brown, \emph{A large values estimate for Dirichlet polynomials}, J. London Math. Soc. (2) 20
  (1979), 8--18.
\bibitem{kerr} B. Kerr, \emph{Large values of Dirichlet polynomials and zero density estimates for the Riemann zeta
  function}, arXiv:1909.12075.
\bibitem{jutila} M. Jutila, \emph{Zero-density estimates for $L$-functions}, Acta Arith. 32 (1977), 55--62.
\bibitem{ingham} A. E. Ingham, \emph{On the estimation of $N(\sigma, T)$}, Quart. J. Math. Oxford Ser. 11 (1940),
  291--292.
\bibitem{montgomery-topics} H. L. Montgomery, \emph{Topics in Multiplicative Number Theory}, Lecture Notes in Math.
  227, Springer, Berlin, 1971.
\bibitem{davenport} H. Davenport, \emph{Multiplicative Number Theory}, 3rd ed., revised by H. L. Montgomery, Graduate
  Texts in Mathematics 74, Springer, New York, 2000.
\bibitem{iwaniec-kowalski} H. Iwaniec and E. Kowalski, \emph{Analytic Number Theory}, American Mathematical Society
  Colloquium Publications 53, American Mathematical Society, Providence, RI, 2004.
\bibitem{friedlander-iwaniec} J. Friedlander and H. Iwaniec, \emph{Opera de Cribro}, American Mathematical Society
  Colloquium Publications 57, American Mathematical Society, Providence, RI, 2010.
\bibitem{tenenbaum} G. Tenenbaum, \emph{Introduction to Analytic and Probabilistic Number Theory}, 3rd ed., Graduate
  Studies in Mathematics 163, American Mathematical Society, Providence, RI, 2015.
\bibitem{huxley72} M. N. Huxley, \emph{On the difference between consecutive primes}, Invent. Math. 15 (1972),
  164--170.
\bibitem{huxley75} M. N. Huxley, \emph{Large values of Dirichlet polynomials, III}, Acta Arith. 26 (1975), 435--444.
\bibitem{halasz} G. Hal\'asz, \emph{\"Uber die Mittelwerte multiplikativer zahlentheoretischer Funktionen}, Acta Math.
  Acad. Sci. Hungar. 19 (1968), 365--403.
\bibitem{montgomery69} H. L. Montgomery, \emph{Mean and large values of Dirichlet polynomials}, Invent. Math. 8
  (1969), 334--345.
\bibitem{montgomery-vaughan} H. L. Montgomery and R. C. Vaughan, \emph{Hilbert's inequality}, J. London Math. Soc.
  (2) 8 (1974), 73--82.
\bibitem{gallagher67} P. X. Gallagher, \emph{The large sieve}, Mathematika 14 (1967), 14--20.
\bibitem{gallagher70} P. X. Gallagher, \emph{A large sieve density estimate near $\sigma = 1$}, Invent. Math. 11
  (1970), 329--339.
\bibitem{kubota} T. Kubota, \emph{On Automorphic Functions and the Reciprocity Law in a Number Field}, Lectures in
  Mathematics, Department of Mathematics, Kyoto University, No.~2, Kinokuniya, Tokyo, 1969.
\bibitem{patterson} S. J. Patterson, \emph{A cubic analogue of the theta series}, J. Reine Angew. Math. 296 (1977),
  125--161; II, ibid., 217--220.
\bibitem{dunn-radziwill} A. Dunn and M. Radziwi\l\l, \emph{Bias in cubic Gauss sums: Patterson's conjecture}, Ann. of
  Math. (2) 200 (2024), 967--1057.
\bibitem{hb-quadratic} D. R. Heath-Brown, \emph{A mean value estimate for real character sums}, Acta Arith. 72 (1995),
  235--275.
\bibitem{goldmakher-louvel} L. Goldmakher and B. Louvel, \emph{A quadratic large sieve inequality over number fields},
  Math. Proc. Cambridge Philos. Soc. 154 (2013), 193--212.
\bibitem{ireland-rosen} K. Ireland and M. Rosen, \emph{A Classical Introduction to Modern Number Theory}, 2nd ed.,
  Graduate Texts in Mathematics 84, Springer, New York, 1990.
\bibitem{harman83} G. Harman, \emph{On the distribution of $\alpha p$ modulo one}, J. London Math. Soc. (2) 27 (1983),
  9--18.
\bibitem{harman} G. Harman, \emph{Prime-Detecting Sieves}, London Math. Soc. Monographs 33, Princeton University
  Press, 2007.
\bibitem{hb-iwaniec} D. R. Heath-Brown and H. Iwaniec, \emph{On the difference between consecutive primes}, Invent.
  Math. 55 (1979), 49--69.
\bibitem{bhp} R. C. Baker, G. Harman and J. Pintz, \emph{The difference between consecutive primes, II}, Proc. London
  Math. Soc. (3) 83 (2001), 532--562.
\bibitem{li-harman} R. Li, \emph{The number of primes in short intervals and numerical calculations for Harman's
  sieve}, arXiv:2308.04458v8 (preprint, 16 October 2025); Theorems~9 and~11 are on p.~36.
\bibitem{banks-shparlinski} W. D. Banks and I. E. Shparlinski, \emph{Bounds on short character sums and $L$-functions
  with characters to a powerful modulus}, J. Anal. Math. 139 (2019), 239--263.
\bibitem{brun} V. Brun, \emph{Le crible d'Eratosth\`ene et le th\'eor\`eme de Goldbach}, Videnskapsselskapets
  Skrifter I, Mat.-naturv. Klasse, Kristiania, 1920, no.~3.
\bibitem{buchstab} A. A. Buchstab, \emph{Asymptotische Absch\"atzung einer allgemeinen zahlentheoretischen
  Funktion}, Rec. Math. [Mat. Sbornik] N.S. 2(44) (1937), 1239--1246.
\bibitem{cheer-goldston} A. Y. Cheer and D. A. Goldston, \emph{A differential delay equation arising from the sieve of
  Eratosthenes}, Math. Comp. 55 (1990), 129--141.
\bibitem{lean4} L. de Moura and S. Ullrich, \emph{The Lean~4 theorem prover and programming language}, Automated
  Deduction -- CADE 28, Lecture Notes in Comput. Sci. 12699, Springer, 2021, 625--635.
\bibitem{mathlib} The mathlib Community, \emph{The Lean mathematical library}, Proceedings of the 9th ACM SIGPLAN
  International Conference on Certified Programs and Proofs (CPP 2020), 367--381.
\bibitem{repo} R. Comley, \emph{number-theory}, the repository of the paper, its supplement and its Lean library,
  \url{https://github.com/dickie81/number-theory}.
